\begin{abstract}
Accurate estimation of vehicle speed is fundamental to autonomous navigation. This study critically evaluates the precision of a vision-based Dense Optical Flow (Farnebäck) algorithm, augmented with dynamic affine calibration, in comparison to consumer-grade GPS trajectory data. Utilizing the comma2k19 dataset, both methodologies were assessed against a 100Hz CAN-bus ground truth within a dynamic urban driving environment. A quantitative approach was employed, involving the calculation of Mean Absolute Error (MAE), Root Mean Square Error (RMSE), and Mean Bias Error (MBE). The findings definitively refuted the initial hypothesis: the GPS-based approach considerably outperformed the vision system, recording an MAE of 1.68 km/h in contrast to 4.50 km/h. Whilst GPS demonstrated considerable robustness across varying environmental conditions, the pure pixel-tracking vision method proved exceptionally vulnerable to visual noise, such as shadows, and failed to reliably detect deceleration. The study concludes that standard GPS probe data is markedly superior to Dense Optical Flow for vehicle speed estimation. Future research should focus on Sensor Fusion techniques, integrating deep learning models (such as YOLO) and Kalman filtering, to address the inherent limitations of both technologies.
\end{abstract}